\section{Introduction}\label{sec:introduction}
Neural-network accelerators at the edge keep their weights in on-chip storage for the whole deployment, so a corrupted weight bit persists across inferences until it is rewritten. Where faults are expected, as in the radiation-hard front-end ASICs of particle detectors, every register is triplicated~\cite{cummings2023econd,diguglielmo2021reconfigurable}. Elsewhere that cost is often unaffordable, which makes selective protection attractive: spend a limited budget on the bits that matter most~\cite{libano2019selective,garciabertoa2023selective,nazari2024fortune,guo2025inpixel}. Almost every proposal does this with a ranking computed once: bit significance, a gradient or a Hessian score~\cite{weng2024fkeras}, or a measured single-flip vulnerability~\cite{tyagi2023thales}. The ranking is then protected from the top down at whatever budget is available.

A fixed ranking assumes that the value of protecting a bit does not depend on what else fails. It does. Protecting a bit removes its effect in a network that already carries the faults of every other unprotected bit, so the value of protection depends on the bit-error probability $p$ and on the rest of the mask, a fact long known in reliability theory~\cite{birnbaum1969importance,rausand2004system}. We found that it matters in practice. Two gradient rankings that differ only in whether they keep the sign of a flip's predicted effect trade places on all nine ResNet-8 checkpoints we studied when both are built from 1,024 calibration inputs: the one that protects loss-increasing flips is ahead by up to 1.6 accuracy points at $p=10^{-3}$ and behind by up to 18 at $10^{-2}$ (Appendix~\ref{app:calibration}). On a language model the same ranking is worse than no protection at all (Appendix~\ref{app:harm}). These observations raise three questions that a designer needs answered. Why does the preferred ranking change with the fault rate? Can the change be predicted before any fault is injected? And can protection be assigned so that it is never far from the best choice, whatever the rate and whatever mechanism provides the protection?

This paper answers them with an expression for the expected loss of a faulty model and a set of tests designed around its predictions.

\paragraph{The expected loss of a faulty model.} Independent stored-bit faults move every weight by a random amount. Expanding the loss to second order around the \emph{average} faulty weights, instead of around the clean ones, splits the expected loss into the loss at those average weights and a curvature-weighted variance (Section~\ref{sec:theory}). Expanded in $p$, the first part contributes one term per bit, which is proportional to $p$ and keeps the sign of a flip's effect, and one term per pair of bits, which is proportional to $p^2$ and penalizes unprotected flips that push the weights in the same direction. A ranking is optimal only while the per-bit term dominates. The pairwise term takes over at a rate $p^*$ that depends on how far the model is from a minimum of its loss and that can be computed from calibration data.

\paragraph{Imperfect protection.} Real mechanisms do not suppress faults; they lower their probability. A triplicated bit is still wrong with probability $3p^2-2p^3$, a bit in a SEC-DED word when another bit of the word also flips, and a weight under checksum correction when another weight of its group is faulty, or when the decoder repairs the wrong weight. Each mechanism therefore has a residual error curve $\rho_s(p)$ (Section~\ref{sec:fault}), and the same expected-loss expression accepts a different residual rate for every bit. Choosing which bits to protect and how to protect them becomes one problem.

\paragraph{Contributions.}
\begin{itemize}
\item \textbf{An explanation of why the preferred ranking depends on the fault rate}, from a second-order expansion of the expected loss around the average faulty weights, and a prediction of the rate $p^*$ at which it switches, computed from calibration data without fault injection (Section~\ref{sec:theory}). \Branch{R1}{The prediction was fixed before the faults were drawn and \TBD{agrees with the measured switch within a factor of R1c} on nine ResNet-8 checkpoints and eight budgets (Section~\ref{sec:switch}).}{It was tested with every prediction fixed before the faults were drawn; \TBD{what held and what did not} (Section~\ref{sec:switch}).}
\item \textbf{Residual error curves for triplication, five copies, SEC-DED codes and checksum correction}, including the wrong repairs of checksum decoders, placed in the same cost (Sections~\ref{sec:fault} and~\ref{sec:theory}). \Branch{R3a}{Each curve matches a bit-level simulation of its mechanism \TBD{R3a}.}{\TBD{which curves did not match simulation}.}
\item \textbf{A rate-aware protection assignment} that minimizes the cost for a given $p$ and storage budget by budget-preserving swaps, with an optional Monte Carlo correction for the large jumps of sign and high-order bits (Section~\ref{sec:method}). \Branch{R2}{On ResNet-8 it is the only tested mask that is never more than half an accuracy point behind the best one at any rate where the model is usable \TBD{largest gap}; the absolute- and positive-gradient rankings each fall behind by \TBD{gaps} somewhere.}{\TBD{Stage 2 outcome, stated as measured}.} \Branch{R3b}{Letting it choose between triplication and five copies per bit improves on ranked triplication by \TBD{R3b}.}{Choosing the mechanism per bit did not improve on ranked triplication \TBD{R3b}.}
\item \textbf{A paired evaluation of seven static rankings on six frozen INT8 models}, from MLPerf Tiny networks to a 0.6-billion-parameter language model, which provides the evidence base (Section~\ref{sec:static} and Appendix~\ref{app:static}). Significance-first protection keeps four of the six models useful only up to about a third of the bit-error probability tolerated by the per-example absolute-gradient mask, and falls 9 to 34 points behind it on every ResNet-8 checkpoint at a moderate fault rate.
\end{itemize}

The study is software-only. Faults are exact XOR flips of stored INT8 weight codes, outcomes come from the quantized executables, and voters and decoders are idealized. We do not measure the area, power or latency of any mechanism, and logical $p$ is not a radiation exposure rate. Our earlier design flow connected susceptibility ranking to selective redundancy and hardware synthesis~\cite{guo2025inpixel}; none of that hardware evidence is reused here.
